%------------------------------------------------------------------------------%
\begin{question}
  Dada a equação da elipse
  \[
    \frac{(x-2)^2}{16} + \frac{(y+1)^2}{9} = 1
  \]
  \vspace{-1.5\baselineskip}
  \begin{enumerate}
    \item identifique o centro e os semieixos ($a$ e $b$)
    \item escreva a equação matricial
    \item esboce a curva
  \end{enumerate}
\end{question}

%------------------------------------------------------------------------------%
\begin{resolution}{Semieixos $a$ e $b$}
  Comparando
  \[
    \frac{(x-2)^2}{16} + \frac{(y+1)^2}{9} = 1
  \]
  \pause
  com a equação canônica transladada
  \[
    \frac{(x-h)^2}{a^2} + \frac{(y-k)^2}{b^2} = 1
  \]
  \pause
  identificamos
  \(
    \qquad
    \pause
    a^2 = 16
    \qquad
    \pause
    a = 4
    \qquad\qquad
    \pause
    b^2 = 9
    \qquad
    \pause
    b = 3
  \)

  \pause
  centro:
  \(
    \qquad
    \pause
    C
    = (h, k)
    \pause
    = (2, -1)
  \)
\end{resolution}

%------------------------------------------------------------------------------%
\begin{resolution}{Equação Matricial}
  Equação Matricial
  \[
    V^t A V = 1
  \]
  \pause
  com
  \[
    A
    = \begin{bmatrix}
      \dfrac{1}{a^2} & 0              \\[2mm]
      0              & \dfrac{1}{b^2}
    \end{bmatrix}
    \pause
    = \begin{bmatrix}
      \dfrac{1}{16} & 0            \\[2mm]
      0             & \dfrac{1}{9}
    \end{bmatrix}
  \]
  \[
    \pause
    V
    = \begin{bmatrix}
      x - h  \\
      y - k
    \end{bmatrix}
    \pause
    = \begin{bmatrix}[c]
      x - 2  \\
      y - (-1)
    \end{bmatrix}
    \pause
    = \begin{bmatrix}
      x - 2  \\
      y + 1
    \end{bmatrix}
  \]
\end{resolution}

%------------------------------------------------------------------------------%
\begin{resolution}{Esboço}
  \centering
  \includegraphics{I-04-Elipse-ex-2}
\end{resolution}

%------------------------------------------------------------------------------%
\endinput
